% Generated draft: requires mathematical and source review.
\documentclass[11pt]{article}
\usepackage[margin=25mm]{geometry}
\usepackage{fontspec}
\setmainfont{Latin Modern Roman}
\usepackage{amsmath,amssymb,mathtools,unicode-math}
\setmathfont{Latin Modern Math}
\usepackage{enumitem,needspace,xurl,hyperref}
\setlength{\parindent}{0pt}
\setlength{\parskip}{5pt}
\setlength{\emergencystretch}{3em}
\newcommand{\N}{\mathbb{N}}
\newcommand{\Z}{\mathbb{Z}}
\newcommand{\Q}{\mathbb{Q}}
\newcommand{\R}{\mathbb{R}}
\newcommand{\C}{\mathbb{C}}
\newcommand{\F}{\mathbb{F}}
\newcommand{\A}{\mathbb{A}}
\newcommand{\PP}{\mathbb{P}}
\newcommand{\E}{\mathbb{E}}
\newcommand{\eps}{\varepsilon}
\newcommand{\dd}{\,\mathrm{d}}
\begin{document}
\section*{UnsolvedMath ID: 30000285}
\textbf{Source problem code:} catalogue entry
\subsection{Definitions and mathematical statement}
For a Riemannian metric $g$ in dimension $n\ge3$, define its Schouten tensor by
\[
 A_g=\frac1{n-2}\left(\operatorname{Ric}_g-
                      \frac{R_g}{2(n-1)}g\right).
\]
For eigenvalues $\lambda=(\lambda_1,\ldots,\lambda_n)$ of $g^{-1}A_g$, let $\sigma_k(\lambda)$ be the $k$th elementary symmetric polynomial. The positive elliptic cone is
\[
 \Gamma_k=\{\lambda:\sigma_j(\lambda)>0\text{ for }1\le j\le k\},
\]
the component of $\{\sigma_k>0\}$ containing the positive orthant. A metric is $k$-admissible if these eigenvalues lie in $\Gamma_k$ everywhere.

The source asks to solve the $\sigma_k$-Yamabe problem on compact manifolds that are not locally conformally flat, in the exact integer range
\[
 2<k\le n/2.
\]
In the usual closed-manifold positive-branch formulation, start with a smooth conformal class containing a $k$-admissible metric and seek a smooth conformal metric $\widetilde g=e^{-2u}g_0$ satisfying
\[
 \lambda(\widetilde g^{-1}A_{\widetilde g})\in\Gamma_k,
 \qquad \sigma_k(\widetilde g^{-1}A_{\widetilde g})=c
 \quad\text{on }M
\]
for some positive constant $c$. The question is to establish this existence in the non-locally-conformally-flat setting throughout the displayed range, including $k=n/2$ when integral.

The terse extract does not state the admissibility or sign branch, whether compact means without boundary, or the exact conformal normalization. These are necessary parts of the conventional formulation above, not a claim that every arbitrary conformal class admits a positive admissible metric. A boundary version would additionally require source boundary conditions and is not supplied.

The literature assessment checked on 21 August 2026 identifies this non-locally-conformally-flat range as unresolved and cites connected-sum work as related evidence. It does not supply a theorem solving all conformal classes in the stated range. This is dated triage requiring independent verification.
\subsection{Short English statement}
Find constant-$\sigma_k$ conformal metrics for $2<k\le n/2$ without assuming local conformal flatness. The standard elliptic problem starts from an admissible conformal class; its branch and closed-manifold conventions must match the source.
\subsection{Sources}
\begin{itemize}
\item \emph{Differentialgeometrie im Grossen} (2006), Oberwolfach Reports, DOI 10.4171/owr/2005/35, pp. 1983--2042, \url{https://doi.org/10.4171/owr/2005/35}; proposed year supplied as 2005.
\item Literature assessment checked 21 August 2026: original report, \url{https://doi.org/10.4171/owr/2005/35}; ``Connected sum construction for sigma-k metrics,'' \url{https://www.researchgate.net/publication/45880998_Connected_sum_construction_for_sigma_k-Yamabe_metrics}.
\item Local UnsolvedMath registry, permanent ID $30000285$, exact source code \texttt{catalogue entry}, created 25 August 2026 and updated 16 September 2026. The clean statement is labelled corrected and verified against the report.
\end{itemize}
\end{document}
